\section{Research Experience}
\cventry{\yearabove{Summer}{2024}--$\,\,$Present}{Graduate Student Researcher}{Applied Nuclear Science Group, University of Michigan}{Ann Arbor, MI}{}{}
\subsection{Particle Identification in YAP:Ce via Machine Learning}
\begin{minipage}[t]{\textwidth}
	\begin{small}
		\begin{itemize}
			\item Compared charge comparison, shape indicator, and random forest (RF) machine learning (ML) methods to discriminate between $\alpha$ and $\gamma$-ray radiation in YAP:Ce.
			\item Quantifed the performance of each method using receiver operating characteristic curves, and demonstrated the substantial advantage of RF ML over traditional methods.
			\item Applied the RF model to a proof-of-concept mixed field radiation environment consisting of $^{241}$Am and $^{137}$Cs sources and demonstrated that known spectral features of $^{241}$Am and $^{137}$Cs are correctly classified.
		\end{itemize}
	\end{small}
\end{minipage}
\subsection{Investigation of the Ionization Yield of Germanium at 254 eV$_{nr}$}
\begin{minipage}[t]{\textwidth}
	\begin{small}
		Collaborator: Dr. Alexander Kavner, University of Zurich
		\begin{itemize}
			\item Jointly conducted an experimental campaign utilizing the spallation neutron source at the Paul Scherrer Institute to re-measure the energy of the final gamma-ray emitted in the $^{72}$Ge(n, $\gamma$)$^{73}$Ge thermal neutron capture reaction, which is necessary to validate a recent ANSG measurement of the ionization yield of Ge.
			\item Developing an algorithm to correct for the non-uniform charge collection of the CdTe detector used in the experiment in order to push the energy resolution closer to the electronic noise + Fano factor limit and therefore produce the lowest possible uncertainty of the final gamma-ray energy.
		\end{itemize}
	\end{small}
\end{minipage}
\subsection{Measurement of the $^{37}$Cl(a, n)$^{40}$K Cross Section}
\begin{minipage}[t]{\textwidth}
	\begin{small}
		Collaborator: Dr. Daniel Santiago-Gonzalez, Argonne National Laboratory
		\begin{itemize}
			\item Analyzing data from the Multi-Sampling Ionization Chamber experiment at ANL, which will validate the stopping power tables for $^{37}$Cl in helium gas and ultimately produce the first measurement of the $^{37}$Cl(a, n)$^{40}$K cross section.
		\end{itemize}
	\end{small}
\end{minipage}
\subsection{}
\cventry{\yearabove{Fall}{2021}--\yearabove{Summer}{2023}}{Undergraduate Student Researcher}{CMS Group, Florida State University}{Tallahassee, FL}{}{}
\subsection{Search for BSM Long-lived Particles}
\begin{minipage}[t]{\textwidth}
	\begin{small}
		Advisor: Prof. Ted Kolberg
		\begin{itemize}
			\item Utilized deep neural network machine learning to search for beyond the Standard Model long-lived particles in data from the CMS Experiment.
			\item Trained ML models on Monte Carlo data with varying signal (mass, lifetime) and background ($t$, $b$ physics) to verify their ability to identify displaced vertices independently of the particular physics.
			\item Investigated the most effective parameters of ML models; verified physics-based expectations.
		\end{itemize}
	\end{small}
\end{minipage}
\subsection{Silicon Sensor Quality Control}
\begin{minipage}[t]{\textwidth}
	\begin{small}
		Advisor: Prof. Rachel Yohay
		\begin{itemize}
			\item Conducted process quality control (PQC) testing for the CMS High Granularity Calorimeter (HGCAL) to monitor the quality and stability of the silicon sensor production process.
			\item Analyzed and presented PQC results to the CMS HGCAL sensor testing group.
			\item Helped create and implement a standardized testing procedure at FSU.
		\end{itemize}

	\end{small}
\end{minipage}
\subsection{}
\cventry{\yearabove{Fall}{2018}--\yearabove{Fall}{2020}}{Undergraduate Student Researcher}{WW Astrochemistry Group, Emory University}{Atlanta, GA}{}{
	Advisor: Prof. Susanna Widicus Weaver
	\begin{itemize}
		\item Performed millimeter/submillimeter spectroscopic measurements of desorbed astrophysical ice analogs.
		\item Conducted proof-of-concept millimeter/submillimeter spectroscopic measurements of H$_2$O and D$_2$O binding energy, sublimation enthalpy, and sublimation entropy.
		\item Analyzed broadband molecular line surveys using the GOBASIC software package.
	\end{itemize}
}
















